\documentclass[11pt,letterpaper]{article}
\usepackage[T1]{fontenc}
\usepackage[utf8]{inputenc}
\usepackage{newtxtext}
\usepackage{amsmath,amssymb,amsthm,mathtools}
\usepackage{newtxmath}
\usepackage[margin=1.02in,headheight=15pt]{geometry}
\usepackage{microtype,booktabs,array,longtable,enumitem,xcolor}
\usepackage{titlesec,fancyhdr}
\usepackage[most]{tcolorbox}
\usepackage{xurl,hyperref,bookmark,needspace}
\definecolor{accent}{HTML}{193F52}
\definecolor{muted}{HTML}{53636C}
\definecolor{pale}{HTML}{F0F5F7}
\hypersetup{colorlinks=true,linkcolor=accent,citecolor=accent,urlcolor=accent,
 pdftitle={Filter Complexity and Booleanization},
 pdfauthor={Research manuscript prepared with ChatGPT},
 pdfsubject={Pseudointersection thresholds, ultrafilter atoms, and Frechet-power sheaves}}
\titleformat{\section}{\Large\bfseries\color{accent}}{\thesection}{0.65em}{}
\titleformat{\subsection}{\large\bfseries\color{accent}}{\thesubsection}{0.65em}{}
\pagestyle{fancy}\fancyhf{}
\fancyhead[L]{\small\color{muted}Filter complexity and Booleanization}
\fancyhead[R]{\small\color{muted}Research manuscript}
\fancyfoot[C]{\small\thepage}
\renewcommand{\headrulewidth}{0.3pt}
\setlength{\parskip}{0.25em}
\setlength{\emergencystretch}{2em}
\setlist{itemsep=0.25em,topsep=0.4em}
\setcounter{tocdepth}{2}
\newtheorem{theorem}{Theorem}[section]
\newtheorem{lemma}[theorem]{Lemma}
\newtheorem{proposition}[theorem]{Proposition}
\newtheorem{corollary}[theorem]{Corollary}
\theoremstyle{definition}
\newtheorem{definition}[theorem]{Definition}
\newtheorem{example}[theorem]{Example}
\theoremstyle{remark}
\newtheorem{remark}[theorem]{Remark}
\newcommand{\om}{\omega}
\newcommand{\Fr}{\mathrm{Fr}}
\newcommand{\Fin}{\mathrm{Fin}}
\newcommand{\B}{\mathbb B}
\newcommand{\C}{\mathbb C}
\newcommand{\PP}{\mathbb P}
\newcommand{\pow}{\mathcal P}
\newcommand{\cp}{\mathfrak p}
\newcommand{\cu}{\mathfrak u}
\newcommand{\cc}{\mathfrak c}
\newcommand{\gen}[1]{\langle #1\rangle}
\newcommand{\Ult}{\operatorname{Ult}}
\newcommand{\RO}{\operatorname{RO}}
\newcommand{\Clop}{\operatorname{Clop}}
\newcommand{\Int}{\operatorname{int}}
\newcommand{\Cl}{\operatorname{cl}}
\newcommand{\Sh}{\operatorname{Sh}}
\newcommand{\Mix}{\operatorname{Mix}}
\newcommand{\Idem}{\operatorname{Idem}}
\newcommand{\At}{\operatorname{At}}
\newcommand{\GE}{\operatorname{GE}}
\newcommand{\supp}{\operatorname{supp}}
\newcommand{\R}{\mathscr R}
\newcommand{\Ssh}{\mathscr S}
\newcommand{\eqset}[2]{E(#1,#2)}
\newcommand{\truth}[1]{\left\lVert #1\right\rVert}
\newcommand{\res}{\mathbin{\upharpoonright}}
\newcommand{\ZFC}{\mathsf{ZFC}}
\newtcolorbox{resultbox}{colback=pale,colframe=accent,boxrule=0.6pt,arc=1mm,
 left=3mm,right=3mm,top=2mm,bottom=2mm,breakable}

\begin{document}
\begin{titlepage}
\thispagestyle{empty}
\vspace*{0.45in}
{\large\sffamily\color{accent}SET THEORY / ORDER THEORY / BOOLEAN-VALUED MODELS\par}
\vspace{0.42in}
{\Huge\bfseries Filter Complexity\par and Booleanization\par}
\vspace{0.23in}
{\LARGE\color{muted}Pseudointersection thresholds, ultrafilter atoms,\par and Fr\'echet-power sheaves\par}
\vspace{0.42in}
{\large A research extension of the ProveIt filter-site report\par}
\vspace{0.13in}
{\large 24 September 2026\par}
\vspace{0.35in}
\begin{resultbox}
For free filters generated by fewer than $\kappa$ sets, the dense-site
Boolean algebra has two exact thresholds:
\[
 \begin{aligned}
 \text{preservation of the original completion}&\quad\Longleftrightarrow\quad\kappa\leq\cp,\\
 \text{absence of atoms and geometric points}&\quad\Longleftrightarrow\quad\kappa\leq\cu.
 \end{aligned}
\]
The first threshold concerns an isomorphism \emph{fixing the original
Boolean events}. Atomicity requires a separate extension property.
An explicit mixing construction computes the associated reduced-power
sheaves and exhibits a sharp infinitary change at $\cp$.
\end{resultbox}
\vfill
\noindent\textbf{Status and scope.} This is an unrefereed research manuscript,
prepared with ChatGPT. The results are proved in ordinary mathematics in
$\ZFC$; they have not been checked in a proof assistant. Classical
cardinal-invariant and Boolean-valued-model machinery is credited rather
than claimed as new. The particular synthesis and applications extend the
inspected repository report, but publication priority has not been established.
One relative-consistency application explicitly imports a published forcing
model; its forcing construction is not reproved here.
\end{titlepage}

\begin{abstract}
The ProveIt report on two-valued gluing obstructions proves that the
Fr\'echet reduced-power presheaf is not a sheaf even on the site of
countably generated free filters. It computes the associated sheaf on the
full filter site as a product of ultrapowers, but that computation cannot
be transferred to the smaller site: its ultrafilters are absent. We compute
the missing Booleanization and place it in a cardinal-indexed hierarchy.
Let $\PP_\kappa$ be the free filters on $\om$ generated, modulo the cofinite
filter, by fewer than $\kappa$ sets. Its Boolean completion is the
regular-open algebra of the topology on $\om^*$ generated by the
ultrafilter-extension sets of these filters. The original
$\pow(\om)/\Fin$ remains order dense exactly for $\kappa\leq\cp$.
Above this threshold there is not even a complete Boolean homomorphism
from its original completion that fixes all original events. Atoms are
exactly ultrafilters of character below $\kappa$, so they first occur at
$\cu^+$. We characterize full atomicity, derive a three-stage example
from Mill\'an's published model, and construct the reduced-power
sheafification by Boolean mixing. Finitary first-order truth is preserved,
whereas a $\cp$-fold conjunction of fixed old predicates separates the
completions sharply. Global rings over fields are von Neumann regular,
with the completion as their idempotent algebra. The countably generated
site is therefore genuinely atomless, not a disguised product of
ultrapowers. A final program distinguishes the questions settled here
from further problems in intermediate atomicity, saturation, choice, and
formalization.
\end{abstract}

\tableofcontents
\newpage

\section{The repository question and the contribution}
\label{sec:scope}

\subsection{A gap left by the existing gluing analysis}
Massas's possibility-semantic construction of the hyperreal line uses
reduced powers indexed by proper filters extending the cofinite filter.
His published question after Fact~5.10 asks whether the uncountability
assumption in his failure-of-sheafness argument is necessary
\cite[p.~517]{Massas}. The ProveIt manuscript
\emph{A two-valued obstruction to gluing Fr\'echet powers} gives a
negative answer using two constants, and also proves failure on the
countably generated subsite \cite{Repo}. We do not claim that resolution
again.

There is, however, an important distinction between the two associated
sheaves. On the full filter site, maximal filters form a dense family;
in $\ZFC$ the associated sheaf has values
\[
 F\longmapsto\prod_{U\supseteq F} M^\om/U.
\]
On the countably generated free-filter site there are no ultrafilter
objects. Reusing the full-site formula on that subsite would give the
wrong answer. Our starting question is thus precise:
\begin{quote}
What is the Booleanization of the smaller filter site, and exactly how
much filter complexity must be admitted before its truth values and its
points change?
\end{quote}
We answer this question and compute the corresponding reduced-power
sheafification. The two-valued obstruction is background motivation,
not an unproved premise of our results.

\subsection{What is classical and what this article adds}
The pseudointersection number $\cp$, the ultrafilter number $\cu$, and
their basic filter interpretations are classical \cite{Blass}. Likewise,
Stone duality, Boolean completion, and the passage from Boolean-valued
models to sheaves by mixing are established tools. The recent treatment
of Pierobon and Viale explicitly develops the relationship between
mixing and dense-topology sheaves \cite{PV}. We give the specialized
proofs needed here so that no semantic convention is left implicit.

The contribution is a single, explicitly marked hierarchy connecting
these tools to the repository's filter site: an exact obstruction to
complete event-preserving comparison, the complete atom and point
spectrum, an early-atomicity application, a concrete sheafification,
and a sharp infinitary truth-value gap. The cardinal thresholds are
translations of classical combinatorics, not newly discovered cardinal
invariants. The resulting statements are research extensions of the
inspected report; their novelty elsewhere is not certified.

The repository snapshot is commit
\texttt{e21766d04c2b8a9b2cdba0cd43e563024b3bd1b9}. The relevant report is
in \path{SetTheory/Cardinals/docs/reports/}
\path{ordinals-and-order-types/frechet-two-valued-obstruction/}.
Its own documentation labels
it unrefereed and not formally verified. We therefore give independent proofs of the repository-specific
structural facts used below instead of inheriting a claim of verification.

\Needspace{18\baselineskip}
\subsection{A map of the principal conclusions}
For an infinite cardinal $\kappa$, let $\C_\kappa$ denote the completion
defined in Section~\ref{sec:representation}. The following table records
the results, with proofs later in the article.
\begin{center}
\small
\begin{tabular}{p{0.52\textwidth}p{0.37\textwidth}}
\toprule
Question & Exact answer\\
\midrule
When does $\C_\kappa$ retain the original completion over
$\pow(\om)/\Fin$? & Exactly when $\kappa\leq\cp$.\\[2pt]
When does $\C_\kappa$ have no atoms or geometric points?
& Exactly when $\kappa\leq\cu$.\\[2pt]
When is it completely atomic?
& Every allowed filter extends an allowed ultrafilter.\\[2pt]
What sheaf replaces $M^\om/F$?
& Arbitrary Boolean mixtures of old sequences.\\[2pt]
When can a zero meet of old events become positive?
& Above $\cp$, first at arity $\cp$.\\[2pt]
What happens for fields?
& A regular global ring with idempotent algebra $\C_\kappa$.\\
\bottomrule
\end{tabular}
\end{center}
Here and throughout, ``over the original Boolean algebra'' means
\emph{fixing its specified embedding}. Abstract isomorphism is a
different question, and Section~\ref{sec:toy} gives an explicit reason
not to confuse them.

\section{Filters, generating character, and dense covers}
\label{sec:filters}

All arguments in this article take place in $\ZFC$. Put
\[
 \Fr=\{A\subseteq\om:\om\setminus A\text{ is finite}\},
 \qquad \B=\pow(\om)/\Fin.
\]
A free filter means a proper filter containing $\Fr$. The notation
$A\subseteq^*D$ means that $A\setminus D$ is finite, and $[A]$ is
the class of $A$ in $\B$. Boolean order on $\B$ is almost inclusion.

\begin{definition}[Generating character modulo finite sets]
For a free filter $F$, set
\[
 \chi_*(F)=\min\{|\mathcal A|:
 F=\gen{\Fr\cup\mathcal A}\}.
\]
For an infinite cardinal $\kappa$ define
\[
 \PP_\kappa=\{F:F\text{ is a free filter and }\chi_*(F)<\kappa\}.
\]
The order is the forcing order: $G\leq F$ means $G\supseteq F$.
Equivalently, the category has an arrow $G\to F$ for every such inclusion.
\end{definition}

The use of $\chi_*$ is deliberate. The cofinite filter has
$\chi_*(\Fr)=0$, although its usual inclusion-base character is
countably infinite. For a nonprincipal ultrafilter the two characters
coincide: closing a family under finite intersections and deleting
finite sets multiplies an infinite size only by $\aleph_0$, while
no nonprincipal ultrafilter has a countable base. The latter fact also
follows from Lemma~\ref{lem:cardinal-bounds} below.

\begin{lemma}[The finite-generation endpoint]
\label{lem:principal}
For every infinite $A\subseteq\om$, let
\[
 F_A=\{D\subseteq\om:A\subseteq^*D\}.
\]
Then $\PP_\om$ consists exactly of the filters $F_A$, and
$[A]\mapsto F_A$ is an order isomorphism from $\B^+$ onto
$\PP_\om$. For any free filter $F$,
\[
 F_A\supseteq F\quad\Longleftrightarrow\quad
 A\subseteq^*D\text{ for every }D\in F.
\]
\end{lemma}
\begin{proof}
A finite generating family can be replaced by its intersection $A$.
Properness implies that $A$ is infinite, and adjoining the cofinite
sets gives exactly $F_A$. Moreover, $F_A\supseteq F_D$ precisely when
$A\subseteq^*D$. The last equivalence is just the definition of $F_A$.
\end{proof}

For $\kappa=\om_1$, the objects are precisely the countably generated
free filters. Every free filter has at most $\cc=2^{\aleph_0}$ generators,
so $\PP_{\cc^+}$ is the full free-filter site.

A sieve $D$ at $F$ is a family of extensions of $F$ closed under further
extension. It is \emph{dense covering} if
\begin{equation}\label{eq:dense-cover}
 \forall G\supseteq F\ (G\in\PP_\kappa\Rightarrow
 \exists H\in D\ (H\supseteq G)).
\end{equation}
The condition is relative to the allowed filters: the same family need
not have the same covering behavior after the site is enlarged.

\begin{lemma}[Small extensions]
\label{lem:small-extensions}
If $F\in\PP_\kappa$ and $A$ is $F$-positive, meaning
$\om\setminus A\notin F$, then $\gen{F\cup\{A\}}$ is a member
of $\PP_\kappa$. The filter generated by two compatible members of
$\PP_\kappa$ is again a member. Thus compatibility can be checked by
ordinary filter compatibility without leaving the site.
\end{lemma}
\begin{proof}
Adjoining $A$ is proper exactly when every member of $F$ meets $A$;
since $F$ contains all cofinite sets, every such intersection is in fact
infinite. The number of generators increases by at most one. Likewise,
a union of two generating families has size below $\kappa$, and its
filter is proper exactly when the original filters are compatible.
Only a finite union of small families is involved, so regularity of
$\kappa$ is not needed.
\end{proof}

\section{A concrete representation of the Booleanization}
\label{sec:representation}

\subsection{A topology on the set of free ultrafilters}
Let $X=\Ult(\B)=\om^*$, with its Stone topology $\tau$. If $b=[A]\in\B$,
write
\[
 \widehat b=\{U\in X:A\in U\}.
\]
These sets are clopen, and they form a base for $\tau$. For a free
filter $F$ put
\[
 K_F=\{U\in X:F\subseteq U\}.
\]
The ultrafilter lemma gives $K_F\ne\varnothing$.

\begin{definition}[The small-intersection topology]
Let $\tau_\kappa$ be the topology on $X$ generated by the sets $K_F$,
$F\in\PP_\kappa$. Equivalently, it is generated by nonempty
intersections of fewer than $\kappa$ original Stone clopens. Put
\[
 \C_\kappa=\RO(X,\tau_\kappa),
 \qquad e_\kappa:\B\longrightarrow\C_\kappa,
 \quad e_\kappa(b)=\widehat b.
\]
We keep this explicit convention instead of using an ambiguous
$G_\kappa$-modification notation.
\end{definition}

\begin{lemma}[Recovery and the clopen base]
\label{lem:recovery}
For free filters $F,G$ and $A\subseteq\om$,
\[
 \begin{split}
 A\in F&\quad\Longleftrightarrow\quad K_F\subseteq\widehat{[A]},\\
 G\supseteq F&\quad\Longleftrightarrow\quad K_G\subseteq K_F.
 \end{split}
\]
The sets $K_F$, $F\in\PP_\kappa$, form a base of nonempty
$\tau_\kappa$-clopens. Two such sets intersect exactly when their
filters are compatible, and a nonempty intersection is the set
corresponding to the generated joint filter.
\end{lemma}
\begin{proof}
The forward implication in the first equivalence is immediate. If
$A\notin F$, then $F\cup\{\om\setminus A\}$ generates a proper
filter: otherwise some member of $F$ would be contained in $A$.
An ultrafilter extension of this filter belongs to $K_F$ but not to
$\widehat{[A]}$. This proves the equivalence and, by applying it to
every member of $F$, proves the second one.

If $\mathcal A$ generates $F$ modulo $\Fr$, then
$K_F=\bigcap_{A\in\mathcal A}\widehat{[A]}$. It is closed in $\tau$
and hence in the finer topology $\tau_\kappa$; it is open by definition.
Finite intersections remain of the same form, by
Lemma~\ref{lem:small-extensions}. The compatibility assertion follows
by extending the joint proper filter to an ultrafilter.
\end{proof}

\subsection{Regular opens and the dense topology}
For a topological space $Y$, an open set $V$ is regular open when
$V=\Int\Cl V$. Its regular-open algebra is a complete Boolean algebra,
with operations
\begin{equation}\label{eq:regular-operations}
 \bigvee_i V_i=\Int\Cl\!\bigcup_iV_i,
 \qquad \bigwedge_i V_i=\Int\!\bigcap_iV_i,
 \qquad \neg V=\Int(Y\setminus V).
\end{equation}
For completeness, if $W=\Int\bigcap_iV_i$ and each $V_i$ is regular
open, then $\Int\Cl W\subseteq\Int\Cl V_i=V_i$ for every $i$.
Being open, $\Int\Cl W$ is therefore contained in $W$; the reverse
inclusion holds for every open $W$. This verifies the meet formula.
The other formulas follow by the same regularization argument.

\begin{theorem}[Representation of the filter-site Booleanization]
\label{thm:representation}
The map $F\mapsto K_F$ identifies $\PP_\kappa$ with an order-dense
subposet of $\C_\kappa^+$. The algebra $\C_\kappa$ is its Boolean
completion. A sieve $D$ covers $F$ exactly when
\begin{equation}\label{eq:cover-join}
 \bigvee_{G\in D}K_G=K_F\qquad\text{in }\C_\kappa.
\end{equation}
Consequently the dense-site sheaf category is
\[
 \Sh(\PP_\kappa,J_{\rm dense})\simeq\Sh(\C_\kappa),
\]
where the right side is the sheaf category of the complete Boolean
algebra viewed as a locale, with covering families having join equal
to the object covered.
\end{theorem}
\begin{proof}
Lemma~\ref{lem:recovery} gives the order embedding. Every nonempty
regular open set contains a nonempty basic clopen $K_F$, proving
order density. Separativity follows as well: if $K_G\nsubseteq K_F$,
the nonempty open difference $K_G\setminus K_F$ contains a condition
incompatible with $F$.

Let $V=\bigcup_{G\in D}K_G\subseteq K_F$. If $D$ is dense in the
sense of \eqref{eq:dense-cover}, every nonempty basic open in $K_F$
contains a member $K_H$ with $H\in D$. Thus $V$ is dense in $K_F$,
which is clopen, and its regularization is $K_F$.
Conversely, suppose $V$ is dense in $K_F$ and take $H\supseteq F$.
Then $K_H$ meets some $K_G$ with $G\in D$. The joint filter is allowed
and belongs to $D$ because $D$ is a sieve. It extends $H$, proving
\eqref{eq:dense-cover}.

Finally, restriction from the Boolean algebra to its dense subposet
preserves covers by \eqref{eq:cover-join}. It also loses no sheaf data:
every nonzero Boolean element is covered by the conditions below it,
and a section over it is uniquely a matching family on that cover.
This describes an inverse extension functor. The same matching-family
construction extends morphisms, and restricting and extending in
either order gives a natural isomorphism. This proves the claimed
equivalence without assuming that the locale has points.
\end{proof}

In particular, $\tau_\om=\tau$ and $\C_\om=\RO(\om^*)$ is the
ordinary order-dense completion of $\B$. At the opposite endpoint,
every ultrafilter is allowed in $\PP_{\cc^+}$, so every singleton
$K_U=\{U\}$ is open. Hence
\begin{equation}\label{eq:full-completion}
 \C_{\cc^+}=\pow(\om^*).
\end{equation}
This power-set algebra is generally \emph{not} the order-dense
completion of the original $\B$.

\section{The exact pseudointersection threshold}
\label{sec:p}

\subsection{The two cardinal invariants}
A family $\mathcal A\subseteq[\om]^\om$ has the strong finite
intersection property, or SFIP, if every finite intersection is infinite.
An infinite $A$ is a pseudointersection of $\mathcal A$ if
$A\subseteq^*D$ for every $D\in\mathcal A$. Define
\[
 \begin{split}
 \cp&=\min\{|\mathcal A|:\mathcal A\text{ has SFIP and no infinite
 pseudointersection}\},\\
 \cu&=\min\{\chi_*(U):U\text{ is a free ultrafilter}\}.
 \end{split}
\]
These are the usual cardinal characteristics \cite{Blass}.

\begin{lemma}\label{lem:cardinal-bounds}
In $\ZFC$, $\aleph_1\leq\cp\leq\cu\leq\cc$.
\end{lemma}
\begin{proof}
For a countable SFIP family $D_0,D_1,\ldots$, recursively choose a
strictly increasing $a_n\in\bigcap_{i\leq n}D_i$. The set of the
$a_n$ is an infinite pseudointersection. Thus $\cp>\aleph_0$.

No free ultrafilter has an infinite pseudointersection. Indeed, if
$A$ were such a pseudointersection, split it into two infinite pieces
$A_0,A_1$. If $A_0\in U$, the inclusion $A\subseteq^*A_0$ fails;
if $A_0\notin U$, then $\om\setminus A_0\in U$, and
$A\subseteq^*(\om\setminus A_0)$ fails. A generating family for
$U$ therefore has no pseudointersection, proving $\cp\leq\cu$.
Finally a free ultrafilter exists, and its size is at most $\cc$.
\end{proof}

\subsection{Preservation and its sharp failure}
\begin{theorem}[Exact marked-completion threshold]
\label{thm:p-threshold}
For every infinite cardinal $\kappa$, the following are equivalent.
\begin{enumerate}[label=\textup{(\roman*)}]
\item $\kappa\leq\cp$.
\item Every $F\in\PP_\kappa$ has an extension of the form $F_A$.
\item $e_\kappa[\B]$ is order dense in $\C_\kappa$.
\item There is a complete Boolean isomorphism
$h:\C_\om\to\C_\kappa$ satisfying $h e_\om=e_\kappa$.
\item There is a complete Boolean homomorphism
$h:\C_\om\to\C_\kappa$ satisfying $h e_\om=e_\kappa$.
\end{enumerate}
When the isomorphism in \textup{(iv)} exists, it is unique.
\end{theorem}
\begin{proof}
Assume $\kappa\leq\cp$, and choose fewer than $\kappa$ generators
for $F$. They have SFIP and hence have an infinite pseudointersection
$A$. Lemma~\ref{lem:principal} gives $F_A\supseteq F$, proving
\textup{(ii)}. Conversely, such an extension is exactly a
pseudointersection of the generators. If $\kappa>\cp$, a family
witnessing $\cp$ generates an allowed filter with no such extension.
Thus \textup{(i)} and \textup{(ii)} are equivalent.

By Theorem~\ref{thm:representation}, every positive Boolean element
contains a $K_F$. Condition \textup{(ii)} says that every such $K_F$
contains an original nonempty clopen $\widehat{[A]}=K_{F_A}$, so
\textup{(ii)} implies \textup{(iii)}. Conversely, if
$\widehat{[A]}\subseteq K_F$, Lemma~\ref{lem:recovery} gives
$F_A\supseteq F$. Hence \textup{(iii)} implies \textup{(ii)}.

An order-dense Boolean embedding has a unique completion up to a
complete isomorphism fixing the embedded algebra. One can see this
directly here: each element of either completion is the join of the
original Boolean elements below it; its set of such lower bounds
is determined by which original Boolean elements are incompatible
with it. The same incompatibility relation holds in both embeddings.
Mapping these cuts to each other gives an order isomorphism and hence
a complete Boolean isomorphism. This proves \textup{(iii)} to
\textup{(iv)}, including uniqueness. The implication from
\textup{(iv)} to \textup{(v)} is immediate.

It remains to show that even \textup{(v)} fails when $\kappa>\cp$.
Let $\mathcal A$ be an SFIP family of size $\cp$ with no infinite
pseudointersection, and let $F=\gen{\Fr\cup\mathcal A}$. The set
\[
 K=K_F=\bigcap_{A\in\mathcal A}\widehat{[A]}
\]
is nonempty and $\tau_\kappa$-clopen. Its original Stone interior is
empty: a nonempty $\widehat{[D]}\subseteq K$ would make $D$ a
pseudointersection. Formula~\eqref{eq:regular-operations} therefore gives
\begin{equation}\label{eq:meet-obstruction}
 \bigwedge_{A\in\mathcal A}e_\om([A])=0,
 \qquad
 \bigwedge_{A\in\mathcal A}e_\kappa([A])=K>0.
\end{equation}
A complete Boolean homomorphism fixing the original events would
preserve these meets, which is impossible.
\end{proof}

\begin{corollary}[The countably generated site]
\label{cor:countable}
The dense Booleanization of the countably generated free-filter site
is canonically $\RO(\om^*)$ over $\pow(\om)/\Fin$. In particular,
it is not $\pow(\om^*)$: the former is atomless and the latter is
atomic. More generally the marked completion is unchanged throughout
$\om\leq\kappa\leq\cp$; its first possible change is at $\cp^+$.
\end{corollary}
\begin{proof}
Apply Theorem~\ref{thm:p-threshold} and
Lemma~\ref{lem:cardinal-bounds}. The algebra $\B$ is atomless,
because every infinite subset of $\om$ splits into two infinite
pieces. Its order-dense completion is consequently atomless as well.
\end{proof}

The obstruction in \eqref{eq:meet-obstruction} is stronger than failure
of density. It rules out \emph{any} complete event-preserving map in
the displayed direction, not just an isomorphism. It does not rule
out an unmarked abstract isomorphism. It also does not assert that the
identity on the underlying set of ultrafilters induces a morphism of
Booleanized locales; the obstruction explains exactly why that
assumption can fail.

\subsection{An exact arity invariant and an explicit inverse map}
\begin{definition}
Let $d_\kappa$ be the least size of a family $\mathcal A\subseteq\B$
whose meet is zero in $\C_\om$ and positive in $\C_\kappa$, with
$d_\kappa=\infty$ if no such family exists. Both meets use the
specified embeddings $e_\om,e_\kappa$.
\end{definition}

\begin{proposition}[The first zero-meet gap]
\label{prop:arity}
\[
 d_\kappa=\begin{cases}
 \infty,&\kappa\leq\cp,\\
 \cp,&\kappa>\cp.
 \end{cases}
\]
\end{proposition}
\begin{proof}
Below the threshold, the complete isomorphism preserves every meet.
Above it, \eqref{eq:meet-obstruction} gives $d_\kappa\leq\cp$.
If a family of fewer than $\cp$ old events has a positive meet in
$\C_\kappa$, all of its finite meets are nonzero in $\B$. Choose
subset representatives; they have SFIP. A pseudointersection gives
a nonzero original Boolean lower bound, so the meet in $\C_\om$
cannot be zero. This gives the reverse inequality.
\end{proof}

\begin{proposition}[How a small filter is represented in the old completion]
\label{prop:interior}
Assume $\kappa\leq\cp$ and $F\in\PP_\kappa$. Then
\[
 K_F=\Cl_\tau\Int_\tau K_F.
\]
Under the inverse isomorphism $\C_\kappa\to\C_\om$ fixing $\B$,
the element $K_F$ is sent to $\Int_\tau K_F$.
\end{proposition}
\begin{proof}
Every original neighborhood of a point of $K_F$ contains a clopen
$\widehat{[D]}$ meeting $K_F$. The filter generated by $F$ and $D$
has fewer than $\cp$ generators, so it has a pseudointersection $A$.
Then $\widehat{[A]}$ lies in $K_F\cap\widehat{[D]}$. Thus the
interior of $K_F$ is dense in $K_F$, proving the first identity.
The interior of a closed set is regular open. It is the join, in the
old completion, of all original clopens contained in $K_F$, which
is exactly the cut description of the inverse isomorphism.
\end{proof}

For the repository's row--tail filter, the countable intersection
$K_{F_\infty}$ is closed but not originally clopen. It nevertheless
has dense original interior, and that interior, not the raw subset
$K_{F_\infty}$, is its representative in $\RO(\om^*)$. Thus the new
calculation agrees with the old gluing obstruction without treating
a closed set as if it were automatically a regular-open event.

\section{The ultrafilter threshold: atoms, points, and atomicity}
\label{sec:atoms}

\subsection{The exact atom spectrum}
Define
\[
 S_\kappa=\{U\in\om^*:\chi_*(U)<\kappa\}.
\]
This is a set of ultrafilters, not a set of hypothetical generic filters
in an extension. The topology $\tau_\kappa$ is Hausdorff because it
contains the original Stone topology.

\begin{theorem}[Atoms are exactly small-character ultrafilters]
\label{thm:atoms}
The atoms of $\C_\kappa$ are precisely the singleton sets
$\{U\}$ with $U\in S_\kappa$. Consequently
\[
 \C_\kappa\text{ is atomless}\quad\Longleftrightarrow\quad
 \kappa\leq\cu.
\]
The first cardinal stage with atoms is $\cu^+$.
\end{theorem}
\begin{proof}
A nonempty regular open set containing two distinct points can be split
by an original Stone clopen separating those points. Its intersections
with that clopen and its complement are nonzero disjoint regular open
sets. Such a set cannot be an atom. Every atom is therefore a singleton.

A singleton $\{U\}$ is open in $\tau_\kappa$ exactly when a basic
$K_F$ contains $U$ and is contained in $\{U\}$. In that case
$K_F=K_U$, and Lemma~\ref{lem:recovery} gives $F=U$; thus
$\chi_*(U)<\kappa$. Conversely, if $U\in S_\kappa$, its singleton
is the basic clopen $K_U$, and hence is an atom. By the definition of
$\cu$, the set $S_\kappa$ is empty exactly for $\kappa\leq\cu$.
\end{proof}

\begin{proposition}[Geometric points of the Booleanization]
\label{prop:points}
The points of the localic topos
$\Sh(\PP_\kappa,J_{\rm dense})$ are indexed by $S_\kappa$.
In particular, it has no points when $\kappa\leq\cu$.
\end{proposition}
\begin{proof}
A point of a locale is a map from its frame to $2$ preserving finite
meets and arbitrary joins. Let $q:\C\to2$ be such a map from a
nontrivial complete Boolean algebra. Put
\[
 z=\bigvee\{b\in\C:q(b)=0\},\qquad a=\neg z.
\]
Join preservation gives $q(z)=0$, so $q(a)=1$ and $a\ne0$.
If $0<b\leq a$, then $q(b)$ cannot be zero, because that would
imply $b\leq z$. If both $b$ and $a\wedge\neg b$ were nonzero,
they would both map to $1$, contradicting preservation of their
zero meet. Thus $a$ is an atom. Conversely an atom $a$ defines
$q_a(b)=1$ exactly when $a\leq b$; completeness and atomicity of
$a$ verify arbitrary-join preservation. These constructions are
inverse. Apply Theorems~\ref{thm:representation} and~\ref{thm:atoms}.
\end{proof}

This does not deny that $X=\om^*$ contains many ordinary points.
An ordinary ultrafilter on a complete Boolean algebra generally
preserves only finite Boolean operations, not arbitrary joins.
Booleanization can therefore have no geometric points even though
the topological presentation from which it was obtained has points.
This distinction is essential in interpreting the countable-site model.

\subsection{Full atomicity is an extension property}
A complete Boolean algebra is atomic if every nonzero element contains
an atom. Having one atom, or many atoms, is weaker.

\begin{theorem}[Atomicity and the canonical decomposition]
\label{thm:atomicity}
For every infinite $\kappa$, the following are equivalent:
\begin{enumerate}[label=\textup{(\roman*)}]
\item $\C_\kappa$ is atomic;
\item $S_\kappa$ is dense in $(X,\tau_\kappa)$;
\item every $F\in\PP_\kappa$ extends to an ultrafilter
$U\in\PP_\kappa$.
\end{enumerate}
Let
\[
 a_\kappa=\bigvee_{U\in S_\kappa}\{U\}
          =\Int_{\tau_\kappa}\Cl_{\tau_\kappa}S_\kappa.
\]
There is a canonical complete Boolean decomposition
\begin{equation}\label{eq:atomic-split}
 \C_\kappa\cong\pow(S_\kappa)\times
       (\C_\kappa\res\neg a_\kappa),
\end{equation}
and the second factor is atomless. It is the trivial one-element
algebra exactly when the equivalent conditions hold.
\end{theorem}
\begin{proof}
Every nonempty regular open contains an atom exactly when it meets
$S_\kappa$, by Theorem~\ref{thm:atoms}. This is equivalent to
density. Since the $K_F$ form a base, density is equivalent to
$K_F\cap S_\kappa\ne\varnothing$ for every allowed $F$, which
is precisely \textup{(iii)}.

For any complete Boolean algebra and any element $a$, the map
$b\mapsto(b\wedge a,b\wedge\neg a)$ gives the product of its two
relative algebras. Below the join of all atoms, an element is uniquely
the join of the atoms below it. Thus the first relative algebra is
the power set of its atoms, here $S_\kappa$. The complementary
relative algebra contains no atom: an atom in it would be an atom
of the original algebra and hence lie below $a_\kappa$, a contradiction.
\end{proof}

Mill\'an's notation $\GE(\mathcal U,\kappa)$ expresses that every
filter with a base of size below $\kappa$ extends a member of a
specified ultrafilter class $\mathcal U$ \cite[Definition~2.1]{Millan}.
For uncountable $\kappa$, the criterion above is exactly
\[
 \GE(S_\kappa,\kappa).
\]
This identifies full atomicity with a familiar but stronger extension
problem, rather than silently deducing it from $\cu<\kappa$.

\begin{remark}[Why ordinary density is insufficient]
If $S_\kappa$ is nonempty, it is dense in the \emph{original} Stone
topology. To see this, transport one small-character free ultrafilter
along a bijection from $\om$ to any prescribed infinite
$A\subseteq\om$. The resulting ultrafilter contains $A$ and has
the same character. But ordinary Stone density says only that the
sets $K_{F_A}$ meet $S_\kappa$; Theorem~\ref{thm:atomicity} requires
this for every small-generated filter. We make no assertion that
existence alone implies, or refutes, this stronger density.
\end{remark}

\subsection{Closure of the presentations}
\begin{proposition}\label{prop:closure}
If $\kappa$ is regular and infinite, then $\PP_\kappa$ is
$<\kappa$-closed: every descending forcing-order chain of length
less than $\kappa$ has a lower bound.
\end{proposition}
\begin{proof}
In filter notation, such a chain is an increasing chain
$(F_\xi)_{\xi<\delta}$, where $\delta<\kappa$. Its union is a
proper filter: any finitely many of its members occur together at
one stage, where their intersection is nonempty. Choose a small
generating family at each stage. Regularity of $\kappa$ ensures
that the union of these $\delta$ families has size below $\kappa$.
It generates the union filter, which is the required lower bound.
The empty chain is bounded by $\Fr$.
\end{proof}

In particular, the countable-base presentation is countably closed.
This assertion concerns a presentation, not a classification of all
its Boolean-algebra invariants. For singular $\kappa$, the union of
fewer than $\kappa$ small bases can have size $\kappa$, so the proof
does not apply. No singular closure conclusion is being smuggled into
the threshold theorems, which did not need regularity.

\section{Early atomicity and genuinely different atomic stages}
\label{sec:model}

The full stage $\C_{\cc^+}=\pow(\om^*)$ is atomic. Is full filter
complexity necessary for atomicity? The answer is no, relative to a
published forcing model. The following is the only imported
relative-consistency input in this article.

\begin{proposition}[Published model input]
\label{prop:model-input}
Mill\'an's proof of Theorem~3.14 \cite[pp.~34--35]{Millan} supplies,
relative to consistency of $\ZFC$, a model with
\[
 \cc=2^{\aleph_1}=\aleph_2,
 \qquad \GE(\mathcal P\!t,\cc),
 \qquad |\mathcal P\!t|=\cc,
\]
where $\mathcal P\!t$ is the class of $P$-point ultrafilters, and every
member of $\mathcal P\!t$ has character $\aleph_1$.
\end{proposition}

Here a $P$-point is a free ultrafilter in which every countable family
has a pseudointersection belonging to the ultrafilter. We only use
the displayed properties, not additional special features of $P$-points.
The cited construction is a countable-support iteration of rational
perfect set forcing. We do not claim to have proved that forcing theorem.

\begin{lemma}[The number of free ultrafilters]
\label{lem:ultra-number}
In $\ZFC$, $|\om^*|=2^\cc$.
\end{lemma}
\begin{proof}
The upper bound follows because an ultrafilter is a subset of
$\pow(\om)$. For the lower bound, use the countably infinite set
\[
 I=\{(n,H):n<\om,\ H\subseteq2^n\}.
\]
For each real $x\in2^\om$, put
$A_x=\{(n,H):x\res n\in H\}$. Given finitely many distinct reals,
their initial segments are distinct for all sufficiently large $n$.
At every such $n$ choose $H$ to include exactly the initial segments
prescribed by an arbitrary binary pattern. Thus every finite Boolean
combination of the $A_x$ with no repeated index is infinite.

For every function $\varepsilon:2^\om\to2$, the family choosing
$A_x$ when $\varepsilon(x)=1$ and its complement otherwise has
SFIP. Together with the cofinite filter it extends to a free
ultrafilter on $I$. Different functions force different ultrafilters.
Choice allows one extension to be selected for each function, giving
$2^\cc$ distinct ultrafilters. Transfer them to $\om$ by a bijection.
\end{proof}

\begin{theorem}[A three-stage hierarchy from the published model]
\label{thm:early-atomicity}
In the model of Proposition~\ref{prop:model-input},
\[
 \cp=\cu=\aleph_1,
 \qquad
 \begin{array}{c|c}
 \text{stage}&\text{Booleanization}\\ \hline
 \kappa=\aleph_1&\C_{\aleph_1}\cong\RO(\om^*)\text{, atomless}\\
 \kappa=\aleph_2&\C_{\aleph_2}\cong\pow(\cc)\\
 \kappa=\aleph_3&\C_{\aleph_3}\cong\pow(2^\cc).
 \end{array}
\]
In particular, the last two complete Boolean algebras are atomic but
are not isomorphic, even without marking the original events.
\end{theorem}
\begin{proof}
The existence of ultrafilters with character $\aleph_1$ gives
$\cu=\aleph_1$, and Lemma~\ref{lem:cardinal-bounds} then gives
$\cp=\aleph_1$. The first line follows from
Corollary~\ref{cor:countable}.

Every $F\in\PP_{\aleph_2}$ has a $P$-point extension, and every
such extension has character $\aleph_1$. Thus the small-character
ultrafilters are dense in this site, and
Theorem~\ref{thm:atomicity} makes $\C_{\aleph_2}$ atomic.
There are at least $\cc$ such ultrafilters, by the stated number
of $P$-points. There are at most $\cc^{\aleph_1}=\cc$ generating
families of size at most $\aleph_1$: use
$\cc=2^{\aleph_1}$ to compute
$(2^{\aleph_1})^{\aleph_1}=2^{\aleph_1}$. Every member of
$S_{\aleph_2}$ is determined by one such family. Therefore
$|S_{\aleph_2}|=\cc$, proving the middle line.

Finally $\aleph_3=\cc^+$, so the last line follows from
\eqref{eq:full-completion} and Lemma~\ref{lem:ultra-number}.
The two atomic algebras have respectively $\cc$ and $2^\cc$
atoms; Cantor's theorem distinguishes their atom sets and hence
the algebras.
\end{proof}

The conclusion settles the proposed ``atomic only at the full site''
possibility negatively, with its exact external assumption exposed.
It does not settle the general atomicity criterion in terms of a
short list of cardinal characteristics. Notice also that two atomic
forcing notions can both add no new sets in the usual generic-extension
sense and still have nonisomorphic complete Boolean algebras and
non-equivalent localic sheaf categories. The atom counts here retain
information that the phrase ``trivial forcing'' discards.

For comparison, under CH we have $\cp=\cu=\cc=\aleph_1$.
Then the countable-base stage is atomless, while the very next cardinal
stage is already the full atomic site. The model above separates
atomicity from the admission of all filters.

\section{A countable toy model and the marking issue}
\label{sec:toy}

The following elementary example isolates the complete-homomorphism
obstruction without requiring uncountable cardinal characteristics.
It is not a claim of a new finite--cofinite algebra phenomenon.
Pierobon and Viale also use that algebra to illustrate delicate
completion behavior \cite[Remark~3.7]{PV}; our purpose is to display
the two filter stages explicitly.

Let $\B_0$ be the Boolean algebra of finite and cofinite subsets of
$\om$. Its Stone space is $Y=\om\cup\{\infty\}$, where each
integer is isolated and neighborhoods of $\infty$ are cofinite.
For an abstract Boolean algebra, repeat the filter construction using
proper filters generated by fewer than $\kappa$ Boolean elements.
At the finite-generation stage, the topology is the original Stone
topology. Since $\om$ is a dense set of isolated points,
\[
 \RO(Y)\cong\pow(\om).
\]
The embedding of $\B_0$ is the ordinary inclusion into $\pow(\om)$.
The isomorphism sends a regular open set to its intersection with
$\om$; its inverse is regularization of that subset in $Y$.

At the countable-generation stage, all proper filters of $\B_0$
are allowed. In particular, the cofinite ultrafilter corresponds to
$\{\infty\}$, which becomes open. The topology is now discrete,
and the completion is $\pow(\om\cup\{\infty\})$. A finite Boolean
element $A$ maps to $A$, while a cofinite one maps to
$A\cup\{\infty\}$.

\begin{proposition}[Abstractly isomorphic, but not compatibly complete]
\label{prop:toy}
The two complete Boolean algebras just described are abstractly
isomorphic. No complete Boolean homomorphism from the first to the
second extends the identity of $\B_0$.
\end{proposition}
\begin{proof}
Their atom sets are both countably infinite, so a bijection of those
sets induces an abstract complete isomorphism. On the other hand,
for $T_n=\{m:m\geq n\}$, the meet of the embedded $T_n$ is empty
in the first algebra and $\{\infty\}$ in the second. A complete
homomorphism extending the original embedding would preserve that
meet, which is impossible. Equivalently, the join of the singletons
$\{n\}$ is $1$ in the first algebra but not in the second.
\end{proof}

Thus ``the completions differ over the original events'' is a genuine
mathematical statement, but it must not be strengthened without proof
to ``the completions are abstractly nonisomorphic.'' The example
supplies a concrete audit test for the wording of
Theorem~\ref{thm:p-threshold}.

\section{Reduced powers and their explicit mixing sheaves}
\label{sec:mixing}

\subsection{The raw presheaf is separated}
Let $M$ be a nonempty set. For a free filter $F$, define
\[
 \R_M(F)=M^\om/{\sim_F},
 \qquad f\sim_Fg\ \Longleftrightarrow\
 \eqset{f}{g}:=\{n:f(n)=g(n)\}\in F.
\]
For $G\supseteq F$, restriction sends $[f]_F$ to $[f]_G$.
These definitions make a presheaf on every $\PP_\kappa$.

\begin{proposition}\label{prop:separated}
The presheaf $\R_M$ is separated for the relative dense topology on
$\PP_\kappa$. Its map into its associated sheaf is injective at
every object.
\end{proposition}
\begin{proof}
If $\eqset{f}{g}\notin F$, adjoining its complement gives an
allowed proper extension $H$ of $F$. At every further extension
of $H$, the two sequences remain inequivalent. A dense cover of
$F$ must reach some such extension. Hence two inequivalent sections
cannot become equal throughout a dense cover. The same fact gives
injectivity of the sheafification unit; it will also follow directly
from the construction below.
\end{proof}

Separatedness is not sheafness. It says that gluing, when it exists,
is unique; it does not supply a representative for every matching
family.

\subsection{Mixtures over a Boolean element}
Fix $\kappa$, abbreviate $\C=\C_\kappa$ and $e=e_\kappa$,
and write $\sigma_F=K_F$. An antichain below $b\in\C$ is a
set of pairwise disjoint nonzero elements, each at most $b$.
It is a partition of $b$ when its join is $b$.

\begin{definition}[A mixture]
A mixture over $b$ is a partition $P\subseteq\C^+$ of $b$,
together with one sequence $f_p\in M^\om$ for each $p\in P$.
We write it as $x=\sum_{p\in P}p\cdot f_p$; the notation denotes
pasting, not numerical addition. Over $0$ the empty partition is
the unique representation.

If $x=\sum p\cdot f_p$ and $y=\sum q\cdot g_q$ are mixtures
on the same support $b$, define their equality value by
\begin{equation}\label{eq:mix-equality}
 \truth{x=y}=\bigvee_{p,q}
       p\wedge q\wedge e([\eqset{f_p}{g_q}]).
\end{equation}
Identify $x$ and $y$ when this value is $b$. Let $\Mix_M(b)$
be the resulting set of equivalence classes.
\end{definition}

All representations form a set: partitions are subsets of the set
$\C$, and labels are chosen from the set $M^\om$. There is no
proper-class collection of names in this construction.

\begin{lemma}[Equality, refinement, and restriction]
\label{lem:mix-well-defined}
Equality at support $b$ is an equivalence relation. Replacing a
partition by a refinement and retaining its labels does not change
the mixture. For $c\leq b$, intersecting each partition element
with $c$ and dropping zeros defines a well-defined restriction
$\Mix_M(b)\to\Mix_M(c)$.
\end{lemma}
\begin{proof}
Reflexivity follows because the diagonal terms of
\eqref{eq:mix-equality} join to $b$; symmetry is immediate.
For transitivity, refine three partitions simultaneously. On a
cell $p\wedge q\wedge r$, the ordinary implication
\[
 \eqset{f_p}{g_q}\cap\eqset{g_q}{h_r}
       \subseteq\eqset{f_p}{h_r}
\]
becomes a Boolean inequality after applying $e$. Taking joins over
all cells gives
$\truth{x=y}\wedge\truth{y=z}\leq\truth{x=z}$.
Thus equality at $b$ is transitive. The same formula shows that
refinement preserves equality values. Intersecting every term with
$c$ multiplies the equality value by $c$, proving that restriction
is independent of representatives. Successive restrictions agree
by associativity of Boolean meet.
\end{proof}

The use of arbitrary joins here is legitimate in a complete Boolean
algebra. Finite meet distributes over arbitrary joins: for fixed
$a$, the operation $x\mapsto a\wedge x$ is left adjoint to
$x\mapsto\neg a\vee x$. No stronger complete-distributivity
assumption is being made.

\begin{lemma}[The mixing sheaf]
\label{lem:mix-sheaf}
The functor $b\mapsto\Mix_M(b)$ is a sheaf on the complete Boolean
algebra $\C$ with its join coverage.
\end{lemma}
\begin{proof}
Suppose $b=\bigvee_i b_i$ and $x_i\in\Mix_M(b_i)$ agree on
all overlaps. Choose a maximal disjoint family of nonzero elements
each below some $b_i$. Its join is $b$: otherwise its nonzero
complement inside $b$ would meet some $b_i$ and permit one more
element. For each cell choose such an $i$, restrict $x_i$ to the
cell, and flatten the resulting partitions. Their union is a
partition of $b$ labelled by sequences, hence a mixture $x$.

On overlap with any $b_j$, each of these pieces agrees with $x_j$
by compatibility of the original family. Formula~\eqref{eq:mix-equality}
therefore gives equality value $b_j$, so $x$ is an amalgam.
If two amalgams exist, their equality value contains every $b_i$
and hence their join $b$. They represent the same element.
\end{proof}

\subsection{The associated-sheaf theorem}
\begin{theorem}[Explicit sheafification on every complexity stage]
\label{thm:sheafification}
The associated sheaf of $\R_M$ on $\PP_\kappa$ is
\begin{equation}\label{eq:associated-sheaf}
 \Ssh_{M,\kappa}(F)=\Mix_M(\sigma_F).
\end{equation}
The unit sends $[f]_F$ to the single-label mixture
$\sigma_F\cdot f$. It is injective. Global sections, meaning
sections at the weakest filter $\Fr$, are exactly $\Mix_M(1)$.
\end{theorem}
\begin{proof}
For two sequences on support $\sigma_F$, their equality value is
$\sigma_F\wedge e([\eqset{f}{g}])$. By
Lemma~\ref{lem:recovery}, this equals $\sigma_F$ exactly when
$\eqset{f}{g}\in F$. Thus the unit is well-defined and injective;
it commutes with restrictions.

The functor in \eqref{eq:associated-sheaf} is a sheaf by
Lemma~\ref{lem:mix-sheaf} and Theorem~\ref{thm:representation}.
It remains to verify its universal property, rather than merely
exhibiting a larger sheaf. Let $x$ be a mixture on $\sigma_F$.
For every nonzero partition cell $p$ of $x$, take all conditions
$\sigma_G\leq p$. These conditions together cover $\sigma_F$
by order density and the fact that the cells join to $\sigma_F$.
On each such $G$, the mixture is represented by the raw sequence
labelling $p$. Thus the unit is locally surjective.

Now let $\alpha:\R_M\to T$ be a morphism into a sheaf $T$.
On the covering conditions just constructed, apply $\alpha$ to
the local raw representatives of $x$. They agree in $T$ on all
overlaps: their equality as mixtures on any overlap is witnessed
locally by conditions contained in their sequence-equality event,
and Lemma~\ref{lem:recovery} says that on these conditions their
raw reduced-power classes are equal. Separatedness of $T$ gives
equality on the whole overlap. Glue in $T$ to define the image of
$x$.

Refining the local representations does not change the glued
section, by uniqueness. The same uniqueness makes the construction
natural in $F$ and independent of the mixture chosen to represent
$x$. It extends $\alpha$, and any extension must have these local
values, so it is unique. This is the sheafification universal property.
Finally $\sigma_{\Fr}=X=1$, giving the description of globals.
\end{proof}

This is a specialized mixification construction, not a new general
sheafification method. The prior framework in \cite{PV} clarifies
why fullness of a Boolean-valued model must not be confused with
its having all mixtures. The explicit form above makes that
distinction usable for the repository's changing filter sites.

\begin{corollary}[The binary and finite-valued cases]
\label{cor:binary}
For $M=2$, there is a canonical Boolean-algebra identification
\[
 \Ssh_{2,\kappa}(F)\cong\C_\kappa\res\sigma_F.
\]
For a finite nonempty set $M$, sections at $F$ correspond exactly
to $M$-indexed Boolean partitions of $\sigma_F$, with zero parts
allowed.
\end{corollary}
\begin{proof}
On a raw sequence $f:\om\to2$, the event for value $1$ is
$e([f^{-1}(1)])$. On a mixture it is the join of the corresponding
piecewise events. This determines the mixture: on that event it
is the constant $1$, and on its complement in $\sigma_F$ it is
$0$. Conversely every $c\leq\sigma_F$ is represented by the
mixture $c\cdot1+(\sigma_F\wedge\neg c)\cdot0$.
For a finite set, the same argument splits each sequence according
to its finitely many values and then joins the pieces carrying each
value. Finite preservation by $e$ makes the resulting events a
partition of the support.
\end{proof}

In particular, the binary associated sheaf on the countably generated
site has global Boolean algebra $\RO(\om^*)$, while the full-site
binary associated sheaf has global Boolean algebra $\pow(\om^*)$.
This computes the missing subsite object rather than merely showing
that a proposed product formula cannot work.

\begin{corollary}[Ultrapowers at the atomic stages]
\label{cor:atomic-products}
If the equivalent conditions of Theorem~\ref{thm:atomicity} hold,
then
\[
 \Ssh_{M,\kappa}(F)\cong
 \prod_{\substack{U\in S_\kappa\\U\supseteq F}}M^\om/U.
\]
In particular, at the full stage this recovers the product of all
ultrapowers in the repository report.
\end{corollary}
\begin{proof}
At the atom $\{U\}$, every mixture has only one nonzero Boolean
piece. Its sequence labels are identified exactly by $U$-equality,
so the value there is $M^\om/U$. Atomicity says that the atoms below
$\sigma_F$ cover it. Distinct atoms have zero overlap, so arbitrary
choices of these local sections glue uniquely. This is precisely
the displayed product.
\end{proof}

\section{First-order transfer and a sharp infinitary separation}
\label{sec:logic}

\subsection{Boolean truth on mixtures}
Let $M$ now be a nonempty, set-sized first-order structure in a
set-sized finitary language. Constants and functions are evaluated
pointwise on sequences and piecewise on mixtures. If finitely many
mixtures have different partitions, intersect those partitions to
obtain a common refinement. A finite common refinement has join $1$,
by repeated distribution of finite meet over arbitrary joins.

For a first-order formula $\varphi(\bar v)$ and a tuple of raw
sequences $\bar f$, write
\[
 V_\varphi(\bar f)=
 \{n<\om:M\models\varphi(\bar f(n))\}.
\]
The intended value on a raw tuple is $e_\kappa([V_\varphi(\bar f)])$.
For mixed tuples, Boolean logical connectives use the operations of
$\C_\kappa$, and quantifiers use joins and meets over the set of
global mixtures. Atomic equality is \eqref{eq:mix-equality}; other
atomic relations use the analogous piecewise formula.

\begin{theorem}[Finitary transfer with an attained existential value]
\label{thm:transfer}
Suppose $\bar x$ is a finite tuple of global mixtures written on a
common partition $(b_i)$ as
$\bar x=\sum_i b_i\cdot\bar f_i$. Then every finitary first-order
formula satisfies
\begin{equation}\label{eq:los-mixing}
 \truth{\varphi(\bar x)}_\kappa
 =\bigvee_i b_i\wedge
       e_\kappa([V_\varphi(\bar f_i)]).
\end{equation}
The value is independent of all representations. For every formula
$\exists v\,\psi(v,\bar x)$, there is a global mixture $y$ such that
\begin{equation}\label{eq:maximum}
 \truth{\exists v\,\psi(v,\bar x)}_\kappa
       =\truth{\psi(y,\bar x)}_\kappa.
\end{equation}
In particular, every sentence true in $M$ has value $1$, and every
sentence false in $M$ has value $0$.
\end{theorem}
\begin{proof}
For terms and atomic formulas, the statement is built into the
piecewise definitions. Equality substitution is valid because, on
any piece where the coordinate sequences agree, the ordinary
atomic statements in $M$ agree. This proves well-definedness at the
atomic level. Refining partitions leaves the displayed join unchanged.

Suppose the formula identity has been proved for $\varphi$.
Complementing a value piecewise on a partition gives
\[
 \neg\bigvee_i(b_i\wedge a_i)
       =\bigvee_i(b_i\wedge\neg a_i).
\]
This follows by intersecting both sides with each $b_i$, which join
to $1$. Since $e_\kappa$ preserves finite Boolean operations,
ordinary negation in $M$ proves the identity for $\neg\varphi$.
The proof for conjunction is the same after common refinement;
the other finitary connectives follow.

For the existential step, suppose the identity holds for $\psi$.
For each $i$ and each $n$, choose $g_i(n)\in M$ satisfying
$\psi(g_i(n),\bar f_i(n))$ whenever any witness exists; use a fixed
default element of $M$ otherwise. This is a set-indexed choice,
permitted in $\ZFC$. Then
\[
 V_\psi(g_i,\bar f_i)=V_{\exists v\,\psi}(\bar f_i).
\]
The mixture $y=\sum_i b_i\cdot g_i$ therefore has, by the induction
hypothesis, exactly the right-hand side of
\eqref{eq:los-mixing} for $\exists v\,\psi$.

For any other candidate mixture $z$, refine its partition with the
$b_i$. On each refined cell, ordinary truth in $M$ gives
$\psi(z(n),\bar f_i(n))\Rightarrow\exists v\,\psi(v,\bar f_i(n))$.
The induction hypothesis bounds its value by that same right-hand
side. Taking the join over all candidates proves both the
existential clause and \eqref{eq:maximum}. Universal quantification
follows by negation. This completes the induction and establishes
representation independence for all formulas.

For a sentence, the set $V_\varphi$ is either $\om$ or
$\varnothing$, independently of the partition. Its value is
therefore $1$ or $0$ as asserted.
\end{proof}

The theorem gives an exact internal, Boolean-valued first-order
model of the theory of $M$. It does \emph{not} say that the set of
global sections, equipped with the relation ``has Boolean value
$1$,'' is an ordinary classical model of every sentence of that
theory. For example, a Boolean disjunction can have value $1$
without either disjunct having value $1$. The field example in
Section~\ref{sec:rings} makes this distinction concrete.

The maximum principle here also does not prove saturation for
infinite types. It produces a witness for one existential
first-order formula. Replacing a single formula by an infinite
family introduces complete meets, which can change with the
filter stage.

\subsection{An explicit infinitary formula detecting the threshold}
\begin{theorem}[A $\cp$-fold conjunction separates the stages]
\label{thm:infinitary}
For every $\kappa>\cp$, there are a structure $M$, a tuple consisting
of one fixed old sequence $f$, and an infinitary formula
$\Psi(v)\in L_{\cp^+,\om}$ such that
\[
 \truth{\Psi(f)}_\om=0,
 \qquad
 \truth{\Psi(f)}_\kappa>0.
\]
Every finitary subformula has its old event as the Boolean value
at both stages. Among conjunctions of fixed old events whose value
changes from zero to positive, the arity $\cp$ is minimal.
\end{theorem}
\begin{proof}
Choose an SFIP family $(A_\alpha)_{\alpha<\cp}$ with no infinite
pseudointersection. Take the underlying set of $M$ to be $\om$
and give the language unary predicates $P_\alpha$, interpreted as
$A_\alpha$. Let $f(n)=n$, and form
\[
 \Psi(v)=\bigwedge_{\alpha<\cp}P_\alpha(v).
\]
Interpret this infinitary conjunction as the complete Boolean meet
of the values of its conjuncts. The individual values are
$e_\lambda([A_\alpha])$ at stage $\lambda$, by
Theorem~\ref{thm:transfer}. Equation~\eqref{eq:meet-obstruction}
therefore gives the two claimed values, the second being
$K_{\gen{\Fr\cup\{A_\alpha:\alpha<\cp\}}}$.
The minimality assertion is Proposition~\ref{prop:arity}.
\end{proof}

There is no contradiction with ordinary truth in $M$. The
intersection of all the $A_\alpha$ is finite, since an infinite
intersection would itself be a pseudointersection. Thus the
pointwise event for the whole infinitary conjunction is zero
modulo finite sets. At the larger stage its \emph{Boolean-valued}
infinitary interpretation is positive. Theorem~\ref{thm:transfer}
was explicitly finitary; applying its formula to this infinitary
conjunction would be invalid.

This distinction also explains the complete-homomorphism obstruction
model-theoretically. A complete comparison preserving all old atomic
predicates would preserve their conjunction. The two values prove
that such a comparison cannot exist. We have not changed the
interpretation of a finitary predicate, and we have not proved
anything about the realization of every type of cardinality $\cp$.

\subsection{The stable range includes the whole sheaf model}
\begin{corollary}\label{cor:model-stable}
When $\kappa\leq\cp$, the unique complete isomorphism
$\C_\om\cong\C_\kappa$ fixing $\B$ transports the mixing model
at one stage to the mixing model at the other, preserving all old
sequences and all Boolean truth values for which the semantics is
defined by complete Boolean operations.
\end{corollary}
\begin{proof}
Apply the isomorphism to every partition coefficient in a mixture,
leaving its sequence labels unchanged. The equality formula,
Boolean operations, and arbitrary joins and meets are preserved.
The inverse isomorphism gives the inverse transport. In particular,
the finitary truth formula and the explicitly defined infinitary
conjunction semantics commute with transport.
\end{proof}

This gives more than an equality of abstract subterminal algebras:
within the stable range the associated models themselves have the
same explicit mixing description. At a small filter $F$, the
support is represented in the original completion by
$\Int_\tau K_F$, as in Proposition~\ref{prop:interior}.

\section{Global rings: regularity and the idempotent spectrum}
\label{sec:rings}

Let $K$ be a field with $0\ne1$, and let
\[
 A_{K,\kappa}=\Ssh_{K,\kappa}(\Fr)
\]
be the ring of global sections. Ring operations are the piecewise
operations from the mixing construction.

\begin{theorem}[Regularity and Boolean idempotents]
\label{thm:rings}
The ring $A_{K,\kappa}$ is a commutative von Neumann regular ring:
for every $x$ there is $y$ with $xyx=x$. Its Boolean algebra of
idempotents is canonically $\C_\kappa$:
\begin{equation}\label{eq:idempotents}
 \Idem(A_{K,\kappa})\cong\C_\kappa.
\end{equation}
Under this identification, primitive nonzero idempotents correspond
exactly to the ultrafilters $U$ with $\chi_*(U)<\kappa$.
\end{theorem}
\begin{proof}
Write $x=\sum_i b_i\cdot f_i$. Define the generalized inverse
sequence
\[
 g_i(n)=\begin{cases}
 0,&f_i(n)=0,\\
 f_i(n)^{-1},&f_i(n)\ne0.
 \end{cases}
\]
The mixture $y=\sum_i b_i\cdot g_i$ satisfies $xyx=x$ pointwise
on every piece, hence globally.

For $b\in\C_\kappa$, the element
$\epsilon_b=b\cdot1+(\neg b)\cdot0$ is idempotent. The map
$b\mapsto\epsilon_b$ is injective because its value-$1$ event is
exactly $b$. Products, Boolean complements, and Boolean joins of
these idempotents are respectively
\[
 \epsilon_b\epsilon_c=\epsilon_{b\wedge c},\qquad
 1-\epsilon_b=\epsilon_{\neg b},\qquad
 \epsilon_b+\epsilon_c-\epsilon_b\epsilon_c
       =\epsilon_{b\vee c}.
\]

Conversely suppose $x^2=x$. In a field, $t^2=t$ is equivalent to
$t=0\vee t=1$. By Theorem~\ref{thm:transfer},
$\truth{x=0}\vee\truth{x=1}=1$. The two events are disjoint,
since $0\ne1$. Therefore $x$ agrees with $0$ on one piece and
with $1$ on the other; the equality formula shows
$x=\epsilon_{\truth{x=1}}$. This proves surjectivity.
An idempotent is primitive exactly when it is nonzero and has no
nonzero strictly smaller idempotent in the Boolean order. Hence
primitive idempotents are atoms, now classified by
Theorem~\ref{thm:atoms}.
\end{proof}

\begin{corollary}[A non-product theorem for the countable site]
\label{cor:not-product}
For every field $K$, if $\kappa\leq\cu$, the global ring
$A_{K,\kappa}$ is not isomorphic to a full direct product of
nonzero fields. In particular, the associated global ring on the
countably generated free-filter site is not a product of
ultrapower fields. At an atomic stage, by contrast,
\[
 A_{K,\kappa}\cong\prod_{U\in S_\kappa}K^\om/U.
\]
\end{corollary}
\begin{proof}
A nonempty product of nonzero fields has primitive idempotents,
one for each coordinate. The ring here has none when
$\kappa\leq\cu$, by Theorem~\ref{thm:rings}. The empty product
is the zero ring and is excluded because $A_{K,\kappa}$ has distinct
constant sections $0,1$. The atomic-stage formula is
Corollary~\ref{cor:atomic-products}.
\end{proof}

This is a statement about a \emph{full direct product}; it does
not deny familiar subdirect-product representations of reduced
rings. Nor does it assert that the global ring is a field in the
atomless case. In fact $\B$ has nontrivial events, so
\eqref{eq:idempotents} supplies nontrivial idempotents at every
stage. For $K=\mathbb R$, the sheaf is internally an ordered field
model with Boolean truth, while its global ring is not a single
hyperreal field.

The ring also remembers the marked-completion obstruction. Any ring
map preserves finite Boolean operations on idempotents. If a ring
comparison between the two stages were required additionally to
preserve all joins of idempotents and to fix the characteristic
sequences of subsets of $\om$, its restriction to idempotents would
be the forbidden complete Boolean map of
Theorem~\ref{thm:p-threshold}. Thus above $\cp$ no comparison with
these combined requirements exists. Ordinary ring homomorphisms
need not preserve arbitrary idempotent joins; without that
requirement this argument gives no prohibition.

\section{General Boolean algebras and intermediate comparisons}
\label{sec:general}

The representation theorem is not peculiar to $\pow(\om)/\Fin$.
Let $B$ be any nontrivial Boolean algebra. For each infinite
$\kappa$, take all proper filters of $B$ generated by fewer than
$\kappa$ elements, and put on $\Ult(B)$ the topology generated by
their sets of ultrafilter extensions. The proof of
Theorem~\ref{thm:representation} applies verbatim. This supplies
a convenient general mechanism whose specialized thresholds can
be computed from the algebra itself.

\begin{definition}
Let $\mathfrak p(B)$ be the least size of a centered family in
$B^+$ with no nonzero common lower bound, and set
$\mathfrak p(B)=\infty$ if no such family exists. Here centered
means that every finite meet is nonzero.
\end{definition}

\begin{proposition}[The abstract threshold mechanism]
\label{prop:general-p}
For the filter hierarchy of $B$, the original $B^+$ is order dense
at stage $\kappa$ exactly when every centered family of size
below $\kappa$ has a nonzero lower bound. Equivalently, this holds
when $\kappa\leq\mathfrak p(B)$, interpreting $\infty$ as larger
than every cardinal. When it fails, a centered family witnessing
the failure gives a zero-to-positive meet and obstructs every
complete homomorphism from the original completion that fixes $B$.
\end{proposition}
\begin{proof}
A principal filter generated by $b>0$ extends the filter generated
by $\mathcal A$ exactly when $b$ is a lower bound for
$\mathcal A$. Thus the first assertion is the order-density
criterion. For a centered family with no positive lower bound,
the intersection of its Stone clopens is nonempty by the
ultrafilter lemma but has empty original interior. Once the
family is small enough to generate an allowed filter, that
intersection is open at the new stage. The two meet calculations
are exactly those of \eqref{eq:meet-obstruction}.
\end{proof}

For $B=\pow(\om)/\Fin$ this invariant is the usual $\cp$.
For the finite--cofinite algebra of Section~\ref{sec:toy}, it is
$\aleph_0$. For a finite Boolean algebra it is $\infty$. The
construction therefore explains both examples within one proof,
without claiming that this abstract reformulation defines an
unprecedented invariant.

\begin{proposition}[A criterion for comparing arbitrary stages]
\label{prop:intermediate}
Let $\kappa\leq\lambda$ be infinite cardinals. The inclusion of
filter posets $\PP_\kappa\subseteq\PP_\lambda$ is order dense
exactly when
\begin{equation}\label{eq:intermediate-criterion}
 \forall F\in\PP_\lambda\ \exists G\in\PP_\kappa\quad G\supseteq F.
\end{equation}
When this holds, their completions are canonically completely
isomorphic in a way respecting every condition from
$\PP_\kappa$, and their reduced-power mixing models agree under
this isomorphism.
\end{proposition}
\begin{proof}
The displayed condition is precisely order density. Compatibility
of conditions of $\PP_\kappa$ is the same in both posets by
Lemma~\ref{lem:small-extensions}; hence a dense inclusion gives
the unique common completion. A raw sequence-equality event is
represented by a condition from the finite-generation stage, so
these events are fixed. Transporting partition coefficients then
identifies the mixing constructions as in
Corollary~\ref{cor:model-stable}.
\end{proof}

The converse in this proposition is asserted for \emph{order density
of the inclusion}, not for the existence of an arbitrary abstract
isomorphism. A single complete algebra can have many inequivalent
presentations and many different embeddings of a common subalgebra.

\section{Proof audit and reproducible finite checks}
\label{sec:audit}

\subsection{Foundational and logical boundaries}
The results separate naturally into three groups. The Stone/filter
representation, the two threshold theorems, the mixing construction,
and the ring consequences are proved here in $\ZFC$. The properties
of the forcing model in Proposition~\ref{prop:model-input} are
imported from Mill\'an; the three-stage consequence is proved from
those properties. The source report supplies the motivating
problem and its previous gluing result, but no theorem in this
article depends on trusting its unformalized proof.

Choice is used at visible points: ultrafilter extension; selection
of bases and representatives; maximal antichain refinements;
pointwise witnesses in a general structure; and cardinal
arithmetic arguments. This article does not extend the source
report's choice-free guarantees to the new theory. Some parts
require less than full choice, but an exact reverse-mathematical
classification is left as a further task.

The most important negative checks are the following.
\begin{enumerate}[label=\textup{A\arabic*.},leftmargin=2.8em]
\item \textbf{The bound is strict in the definition of the site.}
The original completion persists at $\kappa=\cp$; change occurs
at $\cp^+$. Atoms first occur at $\cu^+$, not at $\cu$.
\item \textbf{A marked obstruction is not an abstract classification.}
The no-homomorphism theorem fixes the embedded $\B$.
Proposition~\ref{prop:toy} shows why the qualifier is indispensable.
\item \textbf{An atom is not atomicity.}
The latter requires \eqref{eq:intermediate-criterion} with
ultrafilter targets, as stated in Theorem~\ref{thm:atomicity}.
\item \textbf{A Boolean point is not an ordinary ultrafilter.}
A locale point preserves arbitrary joins. This is why the
countable-site Booleanization is pointless.
\item \textbf{Finitary transfer is not infinitary transfer or saturation.}
The explicit $\cp$-conjunction distinguishes these notions.
\item \textbf{An internal field need not have a field of global sections.}
The idempotent algebra is a direct obstruction to that inference.
\item \textbf{Regularity of $\kappa$ is used only where stated.}
The topology and threshold theorems allow singular cardinals;
the closure proof assumes regularity.
\end{enumerate}

\subsection{What the executable supplement checks}
The archive contains \texttt{verify\_finite.py}, using only exact
integer arithmetic and the Python standard library. It checks
finite Boolean analogues, not the infinite cardinal theorems.
For $B=\pow(n)$, a proper filter is determined by its nonempty
kernel $K\subseteq n$. The program verifies the filter/Stone
order reversal, compatibility, relative dense sieves, and the
resulting atom-supported coverage. It then checks the equality
formula and existential-witness calculation on finite products,
and checks the generalized inverse and all idempotents over
$\mathbb F_3$.

These are tests of the algebraic interfaces used in the proofs.
There is no finite experiment capable of verifying that a family
witnesses $\cp$, that an ultrafilter has a particular uncountable
character, or that the cited forcing model exists. The recorded
JSON output is actual output of the supplied script, not a
substitute for those proofs. The README gives the exact command,
and the separate audit file records the scope of every check.
No Lean or Rocq verification is claimed.


The deterministic run supplied with this manuscript produced:
\begin{center}
\small
\begin{tabular}{lr}
\toprule
Check & Number checked\\
\midrule
Filter order and compatibility pairs & 284\\
Relative filter families & 33,508\\
Sieves among those families & 332\\
Dense sieves among those sieves & 189\\
Candidate Boolean frame-point maps & 65,812\\
Mixture equality comparisons & 507,348\\
Pointwise generalized inverses & 363\\
Idempotent Boolean-operation pairs & 1,364\\
Existential instances / candidate witnesses & 360 / 22,140\\
\bottomrule
\end{tabular}
\end{center}
The script also exhibits a nonzero nonunit in $\mathbb F_3^2$ for
which the internal field disjunction still has full Boolean value.
This small example checks the internal/external distinction, not
any statement about ultrafilters on an infinite set.

\subsection{A formalization path}
A practical formalization should begin with an arbitrary Boolean
algebra and postpone $\om^*$ and cardinal arithmetic. First prove
that proper filters correspond to nonempty Stone closed sets and
that finite unions of small bases preserve compatibility. Next
formalize the dense-site/regular-open representation and the
principal-filter density criterion. The pseudointersection theorem
then becomes an instantiation of a purely order-theoretic lemma.

The mixing construction can be formalized independently over an
abstract complete Boolean algebra with a Boolean-valued sequence
structure. The key finite-support proof obligations are equality
transitivity, finite common refinement, and the existential witness
lemma. Only after these are established should one connect the
locale sheaf construction and the cardinal-indexed sites.

The relative-consistency application should remain a clearly
named external hypothesis until its forcing input has actually
been formalized. A file with an admitted statement is not the same
as a checked derivation of that statement; the theorem dependency
ledger must retain that boundary.

\section{Further research questions and proposed projects}
\label{sec:future}

The questions below are a proposed research program. Except where
a specific published problem has been cited, they are not claims
that the literature has established them as open. The point is to
state exact next targets and the results already available for
attacking them.

\begin{enumerate}[label=\textbf{Q\arabic*.},leftmargin=3em]
\item \textbf{Classify intermediate atomicity.}
Determine the consistency possibilities for
$\GE(S_\kappa,\kappa)$ when $\cu<\kappa\leq\cc$.
Theorem~\ref{thm:atomicity} turns this into an exact extension
question; Theorem~\ref{thm:early-atomicity} gives one positive
intermediate example. Is its truth determined by familiar
cardinal characteristics in any substantial range, or does it
require more detailed ultrafilter-character information?

\item \textbf{Determine mixed atomic and atomless regions.}
For specified $\kappa$, analyze when
$0<a_\kappa<1$ in \eqref{eq:atomic-split}. Equivalently, small
ultrafilters exist but fail to occur above some allowed condition.
Describe the relative completion on the complement of their
regular-open support. The first step is a literature audit of
small-character ultrafilter extension failures, not the unsupported
assumption that $\cu<\kappa$ already decides the issue.

\item \textbf{Compute the full stabilization relation.}
For arbitrary $\kappa\leq\lambda$, decide the extension property
\eqref{eq:intermediate-criterion}. The special case
$\kappa=\om$ is settled here exactly by $\cp$. Other pairs ask
whether complicated filters can always be strengthened to filters
with smaller bases. Determine whether atomicity, once reached,
can be lost at a later stage; no monotonicity assertion follows
from the theorems proved here.

\item \textbf{Separate marked and unmarked isomorphism problems.}
Above $\cp$ the completions cannot be compared by a complete map
fixing all original events. Determine when they can nevertheless
be abstractly isomorphic, and which invariants distinguish them.
Atom counts settle the two atomic stages in our model, while the
finite--cofinite example shows that marking can matter even when
atom counts coincide. Cellularities, distributivity spectra, and
relative densities are natural additional tests.

\item \textbf{Analyze singular stages.}
At singular $\kappa$, compute closure and distributivity of
$\PP_\kappa$ and its completion. The union-of-bases argument
fails at a clearly identified point. One should distinguish the
poset having a particular closure property from it possessing a
dense closed presentation, and from distributivity of the complete
Boolean algebra.

\item \textbf{Develop a genuine type-realization theory.}
Characterize the cardinals $\lambda$ for which relevant
$\lambda$-sized Boolean-valued types over the mixing model are
realized. The maximum principle proved here handles one
existential formula, and the $\cp$-conjunction gives a sharp
fixed-parameter obstruction to preserving infinitary truth.
Neither result alone settles saturation. A useful first target
is a precise theorem for countable types in a countable language,
with the meaning of Boolean finite satisfiability fixed in advance.

\item \textbf{Replace the finite ideal.}
Apply Proposition~\ref{prop:general-p} to
$\pow(I)/\mathcal I$ for other ideals and index sets. Compute both
the centered-family lower-bound threshold and the character
spectrum of quotient ultrafilters. Distinguish a genuinely new
calculation for a particular ideal from merely renaming a known
ideal pseudointersection invariant.

\item \textbf{Remove unnecessary choice.}
Develop a point-free version directly from regular downsets of
the filter poset, avoiding an a priori set of Stone points.
Determine which comparison theorems require the Boolean prime
ideal theorem, which use dependent choice, and which require
stronger choice for representative selection. This would connect
the present $\ZFC$ theory more tightly to the source report's
choice-free gluing obstruction.

\item \textbf{Recover filter complexity from algebraic models.}
The idempotent algebra of $A_{K,\kappa}$ is exactly
$\C_\kappa$. Determine what further information the ring retains
about the marked original sequence algebra, its elementary
ultrapower factors, and the stage parameter. Classify ring maps
that preserve complete idempotent joins separately from ordinary
ring homomorphisms, for which our obstruction is insufficient.

\item \textbf{Build a checked library with explicit assumptions.}
Implement the dependency order in Section~\ref{sec:audit},
beginning with the abstract Boolean statements. A realistic
first milestone is a kernel-checked proof of the principal-filter
density equivalence and the complete-homomorphism obstruction;
a second is the general mixing sheaf. The forcing model should
be a separately audited assumption until its proof is available.
\end{enumerate}

\section{Conclusion}

Restricting the filters is not a harmless implementation choice.
For countably generated filters, the associated binary sheaf has
an atomless regular-open Boolean algebra, and associated field
models have global rings that are not full products of fields.
For all filters, the Boolean algebra is the power set of the
ultrafilter set and the global sections are products of ultrapowers.

The hierarchy between these endpoints has two exact initial
thresholds. The pseudointersection number controls when the
original Boolean events stop being order dense, and even when a
complete event-preserving comparison becomes impossible. The
ultrafilter number controls when atoms and geometric points first
appear. Full atomicity is governed by a separate generic-extension
property. A published model shows that atomicity can arrive before
the full site, and that a later atomic stage can still be a
different complete Boolean algebra.

The mixing calculation ties these order-theoretic facts to
first-order model theory without confusing internal truth with
ordinary truth of global sections. Its sharp infinitary
counterexample identifies exactly what is lost at the first
threshold. These are explicit conclusions about the repository's
filter-site construction; they do not assert a resolution of an
unrelated major set-theoretic conjecture or certify publication
priority.

\begin{thebibliography}{9}

\bibitem{Repo}
Vladimir Reshetnikov, \emph{ProveIt}, repository snapshot
\texttt{e21766d04c2b8a9b2cdba0cd43e563024b3bd1b9}, inspected
24 September 2026. Relevant unrefereed manuscript:
\emph{A two-valued obstruction to gluing Fr\'echet powers},
20 September 2026; especially its countably generated-site discussion
and its full-site sheafification theorem.
\newblock\url{https://github.com/VladimirReshetnikov/ProveIt/tree/e21766d04c2b8a9b2cdba0cd43e563024b3bd1b9/SetTheory/Cardinals/docs/reports/ordinals-and-order-types/frechet-two-valued-obstruction}

\bibitem{Massas}
Guillaume Massas, \emph{A Semi-Constructive Approach to the Hyperreal
Line}, Australasian Journal of Logic \textbf{20}(3) (2023), 490--536.
\newblock DOI: \href{https://doi.org/10.26686/ajl.v20i3.8254}{10.26686/ajl.v20i3.8254}.
The motivating question occurs after Fact~5.10, p.~517.
\newblock\url{https://ojs.victoria.ac.nz/ajl/article/download/8254/7451}

\bibitem{Blass}
Andreas Blass, \emph{Combinatorial Cardinal Characteristics of the
Continuum}, in M.~Foreman and A.~Kanamori (eds.),
\emph{Handbook of Set Theory}, Springer, 2010, 395--489.
Author-hosted draft dated 24 November 2003 was consulted for
background; it is not represented as the final published text.
\newblock\url{https://sites.lsa.umich.edu/ablass/wp-content/uploads/sites/1471/2025/10/hbk.pdf}

\bibitem{PV}
Moreno Pierobon and Matteo Viale, \emph{Boolean valued models,
presheaves, and \'etal\'e spaces}, arXiv:2006.14852,
version~6, 23 July 2026. Used for the established relationship
between Boolean-valued mixing and dense-topology sheaves, and for
the distinction between ordinary and complete Boolean morphisms.
\newblock\url{https://arxiv.org/abs/2006.14852v6}

\bibitem{Millan}
Andr\'es Mill\'an, \emph{Remarks on the Generic Existence of
Ultrafilters on $\omega$}, Topology Proceedings \textbf{34}
(2009), 27--37. Definition~2.1 gives generic existence;
Theorem~3.14 and its proof, pp.~34--35, supply the model used in
Proposition~\ref{prop:model-input}.
\newblock\url{https://topology.nipissingu.ca/tp/reprints/v34/tp34003.pdf}

\end{thebibliography}
\end{document}
